\section{Monte Carlo Design}
\label{chap:monte-carlo-design}

\subsection{Baseline Data-Generating Process}
\label{sec:monte-carlo-design}

The Monte Carlo experiment builds on the high-dimensional simulation design of \textcite{chen_huang_tien_2021}, with \(p=100\) observed controls. The baseline design combines endogenous treatment, high-dimensional controls, and two instruments. For each observation \(i\), the structural disturbances satisfy
\begin{equation}
  \begin{pmatrix}
    u_i \\
    \epsilon_i
  \end{pmatrix}
  \sim
  \mathcal{N}\!\left(
    \begin{pmatrix} 0 \\ 0 \end{pmatrix},
    \begin{pmatrix} 1 & 0.3 \\ 0.3 & 1 \end{pmatrix}
  \right).
  \label{eq:dgp-disturbances}
\end{equation}
so that both disturbances have unit variance and
\(\operatorname{Cov}(u_i,\epsilon_i)=0.3\).

The remaining primitive variables are mutually independent and are also independent of the disturbance pair in Equation~\eqref{eq:dgp-disturbances}:
\begin{equation}
  x_{ji}\overset{\mathrm{ind}}{\sim}\mathcal{N}(0,1),
  \quad j=1,\ldots,100,
  \qquad
  z_{1i},z_{2i},v_{1i},v_{2i}
  \overset{\mathrm{ind}}{\sim}\mathcal{N}(0,1).
  \label{eq:dgp-primitives}
\end{equation}
Observations are independent across \(i\). The observed controls are obtained by applying the standard-normal cumulative distribution function to the latent control variables:
\begin{equation}
  X_{ji}=\Phi(x_{ji}),
  \qquad j=1,\ldots,100,
  \label{eq:dgp-controls}
\end{equation}
Thus \(X_{ji}\sim\mathcal{U}(0,1)\), and the observed control vector is
\(X_i=(X_{1i},\ldots,X_{100i})'\).

The executable design constructs two observed instruments as\footnote{The displayed Monte Carlo equations in \textcite{chen_huang_tien_2021} write the corresponding control terms as $x_j$. The released simulation code instead uses the transformed controls $X_j=\Phi(x_j)$. The simulations in our thesis follow that executable implementation.}
\begin{equation}
  Z_{1i}=z_{1i}+v_{1i}+X_{2i}+X_{3i}+X_{4i},
  \label{eq:dgp-instrument1}
\end{equation}
\begin{equation}
  Z_{2i}=z_{2i}+v_{2i}+X_{7i}+X_{8i}+X_{9i}+X_{10i}.
  \label{eq:dgp-instrument2}
\end{equation}
Let \(Z_i=(Z_{1i},Z_{2i})'\), so the number of instrument components is \(d_z=2\). The excluded variables \(z_{1i}\) and \(z_{2i}\) provide the part of the instruments that enters the treatment index directly. The \(v\) components add independent noise, while the included \(X\) terms generate dependence between the instruments and observed controls. Because these primitives are independent of the structural disturbances, their inclusion does not create the treatment endogeneity described below.

In the baseline design, treatment is generated by
\begin{equation}
  D_i=\Phi(z_{1i}+z_{2i}+\epsilon_i).
  \label{eq:dgp-treatment-baseline}
\end{equation}
The treatment is continuous and lies in \((0,1)\). The structural outcome equation is
\begin{equation}
  Y_i=1+D_i+5\sum_{j=1}^{7}X_{ji}+D_i u_i.
  \label{eq:dgp-outcome}
\end{equation}
The outcome depends directly on \(X_{1i},\ldots,X_{7i}\). Together with the controls entering the instruments, \(X_{1i},\ldots,X_{10i}\) form the ten-control oracle set used by Oracle-GMM.

The source of treatment endogeneity can now be seen directly from Equations~\eqref{eq:dgp-disturbances}, \eqref{eq:dgp-treatment-baseline}, and \eqref{eq:dgp-outcome}. The disturbance \(\epsilon_i\) enters the treatment index, while \(u_i\) enters the structural outcome through the interaction \(D_i u_i\). Because \(u_i\) and \(\epsilon_i\) are positively correlated, realizations of the treatment are statistically related to the unobserved component of the structural treatment effect. Joint normality gives
\[
E[u_i\mid\epsilon_i]=0.3\,\epsilon_i.
\]
Since
\[
D_i=\Phi(z_{1i}+z_{2i}+\epsilon_i)
\]
is increasing in \(\epsilon_i\), this implies
\[
\operatorname{Cov}(u_i,D_i)
=
0.3\,E\!\left[
\epsilon_i
\Phi(z_{1i}+z_{2i}+\epsilon_i)
\right]
>0.
\]
Hence \(E[u_i\mid D_i]\) is not generally zero: treatment is endogenous because it is related to an unobserved determinant of the structural response to treatment.

For a fixed treatment value \(d\in(0,1)\), Equation~\eqref{eq:dgp-outcome} implies the potential outcome
\begin{equation}
  Y_i(d)=1+5\sum_{j=1}^{7}X_{ji}+d(1+u_i).
  \label{eq:dgp-potential-outcome}
\end{equation}
Because \(u_i\) is standard normal and independent of \(X_i\), the conditional structural quantile is
\begin{equation}
  Q_{Y_i(d)\mid X_i}(\tau)
  =1+5\sum_{j=1}^{7}X_{ji}
  +d\left[1+\Phi^{-1}(\tau)\right].
  \label{eq:dgp-structural-quantile}
\end{equation}
where \(\Phi^{-1}\) denotes the standard-normal quantile function. The structural quantile treatment effect is therefore
\begin{equation}
  \alpha_0(\tau)=1+\Phi^{-1}(\tau).
  \label{eq:dgp-alpha-true}
\end{equation}
For the five quantiles considered in the simulations, the corresponding true structural quantile treatment effects are approximately
\[
\begin{gathered}
\alpha_0(0.10)=-0.282,\qquad
\alpha_0(0.25)=0.326,\qquad
\alpha_0(0.50)=1.000,\\
\alpha_0(0.75)=1.674,\qquad
\alpha_0(0.90)=2.282.
\end{gathered}
\]
These values serve as the true treatment effects when evaluating point-estimation accuracy, coverage, and power.

\subsection{Instrument Relevance and the \texorpdfstring{\(\kappa\)}{κ}-Extension}
\label{sec:instrument-relevance-extension}

The baseline design of \textcite{chen_huang_tien_2021} holds instrument relevance fixed. We introduce a single parameter \(\kappa\) that reduces the impact of the excluded variables \(z_{1i}\) and \(z_{2i}\) to treatment. Simply replacing \(z_{1i}+z_{2i}\) with \(\kappa(z_{1i}+z_{2i})\) would also reduce the variance of the latent treatment index. Because treatment is obtained by applying \(\Phi\) to that index, this would change the marginal treatment distribution together with instrument relevance.

To avoid this additional change, the design introduces an independent standard-normal variable \(w_i\) that replaces the variance removed from the excluded-instrument component. Let \(w_i\sim \mathcal{N}(0,1)\) be independent of all previously defined variables. For \(\kappa\in[0,1]\), define
\begin{equation}
  d_i^*(\kappa)
  =\kappa(z_{1i}+z_{2i})
  +\sqrt{2(1-\kappa^2)}\,w_i
  +\epsilon_i,
  \label{eq:dgp-kappa-index}
\end{equation}
\begin{equation}
  D_i(\kappa)=\Phi\!\left(d_i^*(\kappa)\right).
  \label{eq:dgp-treatment-kappa}
\end{equation}
The coefficient \(\sqrt{2(1-\kappa^2)}\) preserves the variance contributed by the excluded part of the treatment index. Appendix~\ref{app:kappa-derivation} derives this coefficient and establishes the distributional properties of the construction. Since
\[
\operatorname{Var}(z_{1i}+z_{2i})=2,
\]
the scaled excluded-instrument component contributes \(2\kappa^2\), while the independent \(w_i\) component contributes \(2(1-\kappa^2)\). Their combined variance therefore remains equal to two.

At \(\kappa=1\), the compensating component disappears and the baseline treatment equation is recovered:
\begin{equation}
  D_i(1)=\Phi(z_{1i}+z_{2i}+\epsilon_i),
  \label{eq:dgp-kappa-baseline}
\end{equation}
As \(\kappa\) falls, the direct contribution of \(z_{1i}\) and \(z_{2i}\) to treatment becomes smaller. At the conceptual limit \(\kappa=0\), these excluded components disappear from the treatment index. The simulations use positive values down to \(\kappa=0.10\), so the weakest design represents very low but nonzero designed relevance rather than a formal case of non-identification.

The variance of the complete latent treatment index remains constant:
\begin{equation}
  \operatorname{Var}\!\left[d_i^*(\kappa)\right]
  =2\kappa^2+2(1-\kappa^2)+1=3.
  \label{eq:dgp-kappa-variance}
\end{equation}
Hence \(d_i^*(\kappa)\sim \mathcal{N}(0,3)\) for every \(\kappa\). Since
\(D_i(\kappa)=\Phi(d_i^*(\kappa))\), the marginal treatment distribution is unchanged across the relevance designs.

Instrument relevance, however, changes directly with \(\kappa\). From the construction of the observed instruments and the latent treatment index,
\begin{equation}
  \operatorname{Cov}\!\left(Z_{1i},d_i^*(\kappa)\right)
  =\operatorname{Cov}\!\left(Z_{2i},d_i^*(\kappa)\right)
  =\kappa.
  \label{eq:dgp-kappa-relevance}
\end{equation}
so the instrument--index covariance declines one-for-one with \(\kappa\). Thus \(\kappa\) provides a transparent design parameter for ordering instrument relevance in this Monte Carlo experiment. It should not be interpreted as a universal measure of IVQR identification strength or as a conventional weak-instrument statistic. The first-stage \(F\)-statistic and partial \(R^2\) reported in Chapter~\ref{chap:results} are therefore used only as descriptive diagnostics of the relevance ordering; no Stock--Yogo or \(F<10\) threshold is imposed on IVQR.

The relationship generating treatment endogeneity is also preserved. Because \(u_i\) is independent of \(z_{1i}\), \(z_{2i}\), and \(w_i\), while \(\operatorname{Cov}(u_i,\epsilon_i)=0.3\),
\begin{equation}
  \operatorname{Cov}\!\left(u_i,d_i^*(\kappa)\right)
  =\operatorname{Cov}(u_i,\epsilon_i)=0.3.
  \label{eq:dgp-kappa-endogeneity}
\end{equation}
Together with Equation~\eqref{eq:dgp-kappa-variance}, this implies that the jointly normal pair
\((u_i,d_i^*(\kappa))\) has the same means, variances, and covariance for every \(\kappa\). Its joint distribution is therefore unchanged across the relevance designs. Because treatment is a fixed transformation of the latent index, the joint distribution of \(u_i\) and \(D_i(\kappa)\) is also preserved.

The extension therefore changes the contribution of the excluded \(z\) components to treatment while keeping the marginal treatment distribution, the structural outcome equation, the structural quantile treatment effects, the control distribution, the observed-instrument construction, the number of moment conditions, and the estimator and inference procedures fixed. Conditional relationships involving treatment and the observed instruments necessarily change because changing that relationship is the purpose of the experiment.

\subsection{Simulation Design and Estimator Comparison}

The completed Monte Carlo experiment uses the design dimensions
\begin{equation}
  \begin{aligned}
    n&\in\{500,1000\},
    &\tau&\in\{0.10,0.25,0.50,0.75,0.90\},\\
    \kappa&\in\{1,0.50,0.25,0.10\},
    &R&=500.
  \end{aligned}
  \label{eq:simulation-design}
\end{equation}
Here $n$ is the Monte Carlo sample size, $\tau$ is the quantile index, $\kappa$ is the instrument-relevance design parameter, and $R=500$ is the number of Monte Carlo replications in each $(n,\tau,\kappa)$ design. The Cartesian product contains $40$ distinct $(n,\tau,\kappa)$ designs, each with $500$ replications.

Within each Monte Carlo replication, the same primitive draws of \(u_i\), \(\epsilon_i\), \(X_i\), \(z_{1i}\), \(z_{2i}\), \(v_{1i}\), \(v_{2i}\), and \(w_i\) are reused across all four \(\kappa\) values. This common-random-number construction improves comparability across the relevance settings because differences across \(\kappa\) are evaluated on matched simulated samples. It does not eliminate Monte Carlo uncertainty.

Oracle-GMM uses the known oracle control set; Full-GMM uses all controls without regularization; and DML-IVQR uses regularized high-dimensional nuisances and an orthogonalization-based score construction, with the implementation-specific residualization constant term described in Section~3.3. The design therefore has $120$ design--estimator cells. Every comparison uses the truth $\alpha_0(\tau)=1+\Phi^{-1}(\tau)$ from Equation~\eqref{eq:dgp-alpha-true}.

All estimators use the same sample within a replication. DML-IVQR uses the BC pivotal penalty and same-sample procedure in Section~3.3. Inference is conducted at 95\%, and power uses the alternatives defined below.

\subsection{Performance Measures and Numerical Inference}
\label{sec:performance-measures}

For a fixed $(n,\tau,\kappa,\text{estimator})$ cell, let $\widehat{\alpha}_r(\tau)$ denote the point estimate in replication $r$, and let $\alpha_0(\tau)$ be the true value from Equation~\eqref{eq:dgp-alpha-true}. Replications are indexed by $r=1,\ldots,R$. For simplicity, we omit the estimator superscript and summarize point-estimation performance using
\begin{equation}
  \begin{aligned}
    \operatorname{Bias}(\tau)
      &=\frac{1}{R}\sum_{r=1}^{R}\left[\alpha_0(\tau)-\widehat{\alpha}_r(\tau)\right], \\
    \operatorname{MAE}(\tau)
      &=\frac{1}{R}\sum_{r=1}^{R}\left|\widehat{\alpha}_r(\tau)-\alpha_0(\tau)\right|, \\
    \operatorname{RMSE}(\tau)
      &=\left\{\frac{1}{R}\sum_{r=1}^{R}\left[\widehat{\alpha}_r(\tau)-\alpha_0(\tau)\right]^2\right\}^{1/2}.
  \end{aligned}
  \label{eq:performance-point}
\end{equation}
Measures are computed separately by estimator and design cell. Positive Bias means estimates lie below the truth on average; MAE records absolute error, and RMSE gives greater weight to large errors.

In the Monte Carlo experiment, the additional subscript $r$ identifies the replication. With $c_{0.95}$ denoting the critical value defined below, define
\begin{equation}
  C_r(\tau)
    =\mathbf{1}\!\left\{W_{n,r}\!\left(\alpha_0(\tau)\right)\leq c_{0.95}\right\},
  \qquad
  \widehat{\operatorname{Coverage}}(\tau)
    =\frac{1}{R}\sum_{r=1}^{R}C_r(\tau).
  \label{eq:performance-coverage}
\end{equation}
Coverage evaluates $W_{n,r}$ directly at $\alpha_0(\tau)$ and is therefore independent of the numerical grid. For DML, the pivotal penalty generates fresh Uniform draws on each evaluation, so a direct truth evaluation and a stored grid-profile evaluation at the same numerical candidate need not use identical draws. So, reported DML coverage is the repeated-sample acceptance frequency of the direct candidate-specific procedure. The empirical size is
\begin{equation}
  \widehat{\operatorname{Size}}(\tau)
    =\frac{1}{R}\sum_{r=1}^{R}
      \mathbf{1}\!\left\{W_{n,r}\!\left(\alpha_0(\tau)\right)>c_{0.95}\right\}
    =1-\widehat{\operatorname{Coverage}}(\tau).
  \label{eq:performance-size}
\end{equation}
The equality follows from the implemented rejection rule.

Power is evaluated at deliberately false treatment-effect values
\begin{equation}
  a_{\Delta}(\tau)=\alpha_0(\tau)+\Delta,
  \qquad
  \Delta\in\{-1.00,-0.50,-0.25,0.25,0.50,1.00\}.
  \label{eq:performance-alternatives}
\end{equation}
For each offset, empirical power is
\begin{equation}
  \widehat{\operatorname{Power}}(\Delta,\tau)
    =\frac{1}{R}\sum_{r=1}^{R}
      \mathbf{1}\!\left\{W_{n,r}\!\left(a_{\Delta}(\tau)\right)>c_{0.95}\right\}.
  \label{eq:performance-power}
\end{equation}
The statistic is evaluated directly at each false value, independently of the confidence-region grid; $\Delta=0$ is size, not power.

The numerical confidence-set calculation proceeds by test inversion \parencite{chernozhukov_hansen_2008}. For fixed $\tau$ and estimator, take a candidate structural effect $a$, re-estimate the required nuisance quantities conditional on it, construct the estimator-specific IVQR score, and compute $W_n(a)$. Under the null that $a$ is the true effect, the asymptotic reference distribution has two degrees of freedom because there are two instrument moment components. Collecting every candidate that passes its own test gives the confidence region. It is not chosen around the point estimate: minimizing $W_n(a)$ and retaining candidates below a critical value are distinct operations.

Conventional Wald inference centered on an estimate typically relies more directly on regular local identification. Direct testing of each hypothesized effect remains useful under weak or partial identification. This robustness concerns test validity, not narrow regions: many competing effects may survive under weak relevance.

The accepted-set measure totals all accepted components inside the primary computational domain, including disconnected components. It is neither the span between extreme accepted values nor an unrestricted confidence-region length. Because $d_z=2$, the 95\% cutoff for the statistic in Section~3.4 is
\begin{equation}
  c_{0.95}=\chi^2_{2,0.95}=5.991464547\ldots.
  \label{eq:wn-critical-value}
\end{equation}
A candidate is accepted when $W_n(a)\leq c_{0.95}$ and rejected when $W_n(a)>c_{0.95}$, with the 95\% cutoff $5.991464547\ldots$ above. The primary grid is
\begin{equation}
  \begin{aligned}
    A_0&=[-1,3],
    &h&=0.1,\\
    G_0&=\{a_j=-1+(j-1)h:j=1,\ldots,41\},\\
       &=\{-1,-0.9,\ldots,2.9,3\}.
  \end{aligned}
  \label{eq:primary-grid}
\end{equation}
Here $A_0=[-1,3]$ is the fixed primary computational domain, $h=0.1$ is the distance between adjacent grid points, $j$ is the grid-point index, and $a_j$ denotes the $j$-th candidate treatment-effect value. The set $G_0$ contains all 41 candidate values. There are 41 points because the four-unit interval from $-1$ to $3$, evaluated with a step size of $0.1$ and including both endpoints, contains $4/0.1+1=41$ grid points. The original authors' main Monte Carlo implementation uses \texttt{seq(-1,3,length=41)}. Preserving its four-unit domain and $0.10$ resolution avoids choosing a new grid after observing the extension results. All five quantiles, strength settings, and estimators use this common grid, making their measures comparable without differences in grid density. The resolution is inherited, not claimed to be asymptotically optimal. All five true effects lie strictly inside the domain: the extreme values are $\alpha_0(0.10)\approx-0.282$ and $\alpha_0(0.90)\approx2.282$, with $\alpha_0(0.50)=1$.

The distinction between the theoretical region and its finite-grid implementation is important. The theoretical region $\mathcal C=\{a:W_n(a)\leq c_{0.95}\}$ ranges over the relevant parameter space. Computation evaluates only finitely many candidates, producing an accepted set inside $A_0$, a computational search region rather than an asserted parameter space. The reported summary is therefore the grid-based accepted-set measure within the primary computational domain $A_0=[-1,3]$, rather than unrestricted confidence-region length.

Total confidence-region length is an intuitive informativeness summary, but numerical measurement requires localization inside a sufficiently wide search domain. The initial design considered adaptive expansion; frozen diagnostics showed that weak-relevance accepted sets could keep reaching wider boundaries. Stopping at an arbitrary endpoint and reporting that finite span would make length partly an endpoint artifact. The main comparison therefore uses fixed $A_0$, while separate $A_1=[-3,5]$ and $A_2=[-5,7]$ diagnostics examine endpoint sensitivity (see Appendix A).

For replication $r$, define the binary grid-acceptance indicator, accepted grid, and all-grid-points indicator by
\begin{equation}
  \begin{aligned}
    I_r(a_j)
      &=\mathbf{1}\!\left\{W_{n,r}(a_j)\leq c_{0.95}\right\}, \\
    \mathcal{C}_{r,G_0}
      &=\left\{a_j\in G_0:I_r(a_j)=1\right\}, \\
    F_r
      &=\mathbf{1}\!\left\{I_r(a_j)=1\ \text{for every }a_j\in G_0\right\}.
  \end{aligned}
  \label{eq:grid-acceptance}
\end{equation}
For the primary domain $A_0$, define left-boundary, right-boundary, and either-boundary acceptance in replication $r$ by
\begin{equation}
  B_r^{L}=I_r(a_1),\qquad
  B_r^{R}=I_r(a_{41}),\qquad
  B_r^{E}=\mathbf{1}\{B_r^{L}=1\ \text{or}\ B_r^{R}=1\}.
  \label{eq:boundary-acceptance}
\end{equation}
The corresponding either-boundary acceptance rate over $R$ Monte Carlo replications is
\begin{equation}
  \widehat B^{E}
  =\frac{1}{R}\sum_{r=1}^{R}B_r^{E}.
  \label{eq:boundary-acceptance-rate}
\end{equation}
The same definition is used for an expanded diagnostic domain, with its own left and right endpoints. A boundary is accepted when the candidate-specific test does not reject the endpoint of the finite numerical search interval. A high boundary-acceptance rate therefore means that accepted components frequently reach the edge of the examined domain and are not localized inside it. An accepted endpoint suggests that a component may continue outside the domain; a measure near four indicates limited localization within $A_0$. The $A_1$ and $A_2$ diagnostics examine whether this behavior disappears. Failure to localize even within $A_2$ does not establish global unboundedness or failed identification. If boundary acceptance disappears after the domain is widened, the observed accepted grid components no longer touch the investigated endpoints. This does not rule out additional disconnected accepted components outside the investigated domain. Persistent boundary acceptance as successively wider finite domains are examined indicates continued lack of numerical localization over those investigated ranges.

The frozen implementation computes the grid-based accepted-set measure within $A_0$ as integral of this binary indicator:
\begin{equation}
  M_r(A_0)
  =\sum_{j=1}^{40}\frac{a_{j+1}-a_j}{2}
    \left[I_r(a_j)+I_r(a_{j+1})\right].
  \label{eq:accepted-measure}
\end{equation}
The sum runs from $j=1$ to $40$ because the 41 grid points in $G_0$ define 40 adjacent grid intervals $[a_j,a_{j+1}]$. The sum includes separated accepted components without filling rejected gaps and does not interpolate critical-value crossings. With $h=0.10$, an adjacent interval contributes zero if both endpoints are rejected, $0.10$ if both are accepted, and $0.05$ if exactly one is accepted. Thus the measure ranges from zero to four and summarizes numerical acceptance, including disconnected components, without solving for exact continuous crossings of $W_n(a)$ with the cutoff. The Monte Carlo mean of $F_r$ is the frequency with which all tested grid points are accepted, not acceptance of the real line.

Coverage evaluates the exact $\alpha_0(\tau)$, whether or not it is a grid point; power evaluates the exact $\alpha_0(\tau)+\Delta$, independently of $G_0$. Grid and domain choices affect the accepted-set measure and profile point estimate, but do not mechanically determine reported coverage or power. A post-main diagnostic examines sensitivity to the finite endpoints using
\begin{equation}
  A_0=[-1,3],\qquad A_1=[-3,5],\qquad A_2=[-5,7].
  \label{eq:grid-diagnostic-domains}
\end{equation}
It uses $n=1000$, saved primitive samples, all five quantiles and estimators, $\kappa\in\{1,0.25,0.10\}$, and spacing $0.1$. Stage~1 evaluates $A_1$ and its nested $A_0$ profile; Stage~2 reloads the same samples and adds only the tails needed for $A_2$. Domain comparisons therefore change candidates, not samples. The outputs record accepted measure, grid share and count, boundary acceptance, all-grid acceptance, and incremental measure.

The boundary-acceptance diagnostic follows the finite-domain interpretation of Equations~\eqref{eq:boundary-acceptance}--\eqref{eq:boundary-acceptance-rate}. The bounded diagnostic ends at $A_2$, while main summaries retain $A_0$ and direct coverage and power remain unchanged.
